%=============================================================================
\chapter{Summary and Conclusions}\label{chap:conclusion}
%=============================================================================

\section{Summary of Path Planning Approaches}\label{sec:summary}

Path planning is a vital component of swarm operations that fundamentally affects
mission success, energy consumption, and the efficiency of the whole system.
Advancing these algorithms---often driven by biological inspiration and modern AI
techniques---is essential for realizing the full potential of autonomous systems.

\subsection{Overview of Deterministic, AI-Based, and Hybrid Approaches}

According to their core methodology, the path planning techniques developed for
swarm systems fall into three main categories: deterministic approaches, AI-based
approaches, and hybrid approaches.

\paragraph{Deterministic approaches.} These are classical algorithms that rely on
established mathematical models and predefined rules to generate paths. Although
these methods often provide globally optimal paths and a straightforward
implementation, they face serious difficulties in dynamic and unpredictable
environments and can be computationally expensive; hence, on their own, they are
not commonly used for real-time applications.

\paragraph{AI-based approaches.} These methods use machine learning and other
intelligent technologies to adapt to dynamic and uncertain environments by
learning from data or experience \cite{abujabal2024path}. Their key advantage is the ability to
handle complex scenarios, to make intelligent decisions in real time, and to adapt
quickly to unexpected obstacles or changes in the environment. However, they
generally require extensive training data and substantial computational
resources.

\paragraph{Hybrid approaches.} These strategies integrate elements of both
deterministic and AI-based methods in order to exploit their complementary
strengths and compensate for the limitations of each. The goal is to produce
solutions that are more robust, adaptive, and efficient in complex and dynamic
environments. Studies show that hybrid approaches make up roughly 25\% of the
reviewed path planning literature, which indicates an open research space for
advancing the field.

\subsection{Methodologies and Their Principles}

The core principles of these three families of approaches reveal their distinct
contributions to solving the path planning problem in swarm robotics.

\paragraph{Principles of deterministic methods.} Deterministic methods usually
transform the continuous motion planning problem into a discrete search problem.
\begin{itemize}
  \item \lead{Graph-based methods (e.g., \Astar)} These methods determine the
        optimal path on a discrete map using heuristic cost functions. A common
        formulation of \Astar{} uses the function $f(n) = g(n) + h(n)$, where
        $g(n)$ is the accumulated cost and $h(n)$ is the heuristic estimate of the
        cost to the goal.
  \item \lead{Sampling-based methods (e.g., PRM, RRT, \RRTstar)} These methods
        explore complex, high-dimensional spaces through random sampling. They
        are usually computationally efficient and probabilistically complete,
        but they do not guarantee global optimality.
  \item \lead{Artificial potential fields (APF)} These methods generate the path
        from geometric principles, using a virtual force that attracts the robot
        toward the goal and repels it from the obstacles. Conceptually, the total
        potential can be written as $U(q) = U_{\att}(q) + U_{\rep}(q)$ and the
        guiding force as $F(q) = -\nabla U(q)$, so that the path is shaped by
        gradient descent.
\end{itemize}

\paragraph{Principles of AI-based methods.} These methods rely on learning and on
iterative optimization inspired by natural processes; in path planning,
bio-inspired techniques have been the most widely used.
\begin{itemize}
  \item \lead{Particle swarm optimization (PSO)} A stochastic algorithm based on
        social behavior that balances the exploitation of the best local position
        with the exploration of the best global position. The standard
        formulation consists of the velocity and position updates
        \begin{equation}
          \begin{aligned}
            v_i^{t+1} &= \omega\, v_i^{t}
                         + c_1 r_1\left(\mathit{pbest}_i - x_i^{t}\right)
                         + c_2 r_2\left(\mathit{gbest} - x_i^{t}\right), \\
            x_i^{t+1} &= x_i^{t} + v_i^{t+1},
          \end{aligned}
          \label{eq:pso-summary}
        \end{equation}
        which make clear the roles of the inertia weight $\omega$, the
        acceleration coefficients $c_1$ and $c_2$, and the random factors $r_1$
        and $r_2$ in controlling the balance between exploration and
        exploitation.
  \item \lead{Ant colony optimization (ACO)} Agents mark and follow desirable
        paths using an artificial pheromone, which reduces redundant motion and
        improves coordination. The common evaporation and reinforcement rule is
        $\tau_{ij} \leftarrow (1-\rho)\,\tau_{ij} + \sum_{k} \Delta\tau_{ij}^{k}$
        with $\Delta\tau_{ij}^{k} = Q/L_k$ on the traversed edges, which
        facilitates convergence to shorter paths.
  \item \lead{Reinforcement learning (RL)} Agents learn an optimal control policy
        through trial-and-error interaction and by receiving rewards or
        penalties, which makes RL suitable for dynamic and uncertain
        environments.
\end{itemize}

\paragraph{Principles of hybrid methods.} Hybrid methods integrate the efficiency
of one algorithm with the specialized capability of another. For example,
combining a deterministic algorithm with an AI technique can be used to tune the
parameters adaptively or to guide the search process. A prominent example is the
enhancement of a fuzzy neural network (FNN) path planner by integrating the
\Astar{} algorithm to guarantee path optimality, which improves the rationality
and feasibility of the path by removing limitations such as motion jumps.

%-----------------------------------------------------------------------------
\section{Conclusion}\label{sec:conclusion}

The systematic review of path planning methodologies and real-world applications
in this seminar confirms the vital position of swarm robotics as a key field in
the advancement of autonomous systems. Swarm systems are designed to overcome the
inherent limitations of single-robot or traditional multi-robot configurations,
and they draw on the principles of spontaneous self-organization in nature and on
computational intelligence.

\subsection{Advantages of Swarm Robotics}

\begin{itemize}
  \item \lead{Fault tolerance} Swarm systems are inherently robust owing to their
        extensive redundancy: the failure of one robot does not undermine the
        integrity of the mission. By distributing control and functionality
        across many agents and eliminating single points of failure, high
        robustness and reliability are achieved.
  \item \lead{Scalability} A central feature of swarm robotics is efficient
        scaling from a few robots to thousands while maintaining performance and
        efficiency. This scalability stems directly from local, simple, and
        distributed coordination mechanisms.
  \item \lead{Productivity and parallelism} Swarms have an advantage in complex,
        large-scale problems, performing tasks that are impractical for a single
        robot. Their algorithms are generally designed to be parallel, so that
        multiple goals can be handled simultaneously, yielding significant time
        savings in search, exploration, and construction.
  \item \lead{Decentralization and simplicity} Decentralized control and
        spontaneous self-organization produce macroscopic behaviors from simple
        local interactions. The simpler components of individual robots make
        development easier and cheaper than that of complex centralized units.
\end{itemize}

\subsection{The Critical Role of Advanced Path Planning}

\begin{itemize}
  \item \lead{The need for safe and reliable motion} Path planning is vital for
        ensuring safe, reliable, and smooth motion, especially in large-scale or
        obstacle-rich environments. It directly affects mission success, resource
        utilization, and energy consumption.
  \item \lead{Leveraging collective intelligence} Advanced path planning based on
        swarm intelligence principles, such as ACO, PSO, and RL, is essential for
        realizing the potential of swarm systems. By extracting insight from
        experience and data, these methods provide intelligent adaptation in
        dynamic and uncertain environments---a capability that traditional
        deterministic methods alone lack.
  \item \lead{Hybridization for robustness and efficiency} Developing hybrid
        approaches is key: the optimality of deterministic methods is combined
        with the flexibility of AI-based methods \cite{abujabal2024path}. A prominent example is
        the integration of a fuzzy neural network (FNN) with the \Astar{}
        algorithm, which guarantees path optimality and improves feasibility and
        controllability by removing the inherent limitations of the FNN.
  \item \lead{Scalability through virtualization and optimization} Advanced
        methods must cope with the severe computational complexity of path
        planning for large swarms. The virtual tube optimization method, by
        combining a centralized definition of the path with distributed model
        predictive control (MPC), drastically reduces the computational cost
        through simple interpolation of the optimal path and preserves
        scalability in obstacle-rich environments.
\end{itemize}
